\BourbakiExerciseGroup

\begin{enumerate}
\def\labelenumi{\arabic{enumi})}
\tightlist
\item
  \begin{enumerate}
  \def\labelenumii{\alph{enumii})}
  \tightlist
  \item
    Give an example of continuous mappings \(f: \mathbf{X} \to \mathbf{Y}\) , \(g: \mathbf{Y} \to \mathbf{Z}\) such that \(g \circ f\) is proper, \(f\) is not surjective and \(g\) is not proper.
  \end{enumerate}
\end{enumerate}

\begin{enumerate}
\def\labelenumi{\alph{enumi})}
\setcounter{enumi}{1}
\tightlist
\item
  Give an example of two proper mappings \(f: \mathbf{X} \to \mathbf{Y}\) and \(g: \mathbf{X} \to \mathbf{Z}\) where \(\mathbf{X}, \mathbf{Y}, \mathbf{Z}\) are quasi-compact, such that \(x \to (f(x), g(x))\) is not a proper mapping of \(\mathbf{X}\) into \(\mathbf{Y} \times \mathbf{Z}\) .
\end{enumerate}

\begin{enumerate}
\def\labelenumi{\arabic{enumi})}
\setcounter{enumi}{1}
\tightlist
\item
  Let \(f: \mathbf{X} \to \mathbf{Y}\) be a proper mapping and let \(\mathbf{A}\) be a subspace of \(\mathbf{X}\) . Show that if \(\mathbf{A} = \overline{f}^1(f(\mathbf{A})) \cap \overline{\mathbf{A}}\) then the restriction \(f|_{\mathbf{A}}\) is a proper mapping of \(\mathbf{A}\) into \(f(\mathbf{A})\) . Is this condition necessary?
\end{enumerate}

¶ 3) Let \((\mathbf{X}_{\alpha}, f_{\alpha\beta})\) be an inverse system of quasi-compact Kolmogoroff spaces such that the \(f_{\alpha\beta}\) are proper mappings. Show that if the \(\mathbf{X}_{\alpha}\) are non-empty then \(\varprojlim \mathbf{X}_{\alpha}\) is not empty (show that in a quasi-compact Kolmogoroff space, every non-empty closed subset contains a closed subset consisting of a single point).

\begin{enumerate}
\def\labelenumi{\arabic{enumi})}
\setcounter{enumi}{3}
\item
  Give an example of a continuous mapping \(f: \mathbf{X} \to \mathbf{Y}\) (where \(\mathbf{X}\) and \(\mathbf{Y}\) are Hausdorff spaces) such that the inverse image \(\overline{f}^1(\mathbf{K})\) of every compact subset \(\mathbf{K}\) of \(\mathbf{Y}\) is compact, but which is not proper (cf.~§ 9, Exercise 4).
\item
  Let \(f: \mathbf{X} \to \mathbf{Y}\) be a proper mapping.
\end{enumerate}

\begin{enumerate}
\def\labelenumi{\alph{enumi})}
\tightlist
\item
  Show that if X is regular then \(f(\mathbf{X})\) is a regular subspace of Y.
\item
  Show that if \(f(\mathbf{X})\) is a Lindelöf subspace of Y (§ 9, Exercise 14), then X is Lindelöf (use Proposition 10 of § 5, no. 4).
\end{enumerate}

¶ 6) Let X be a space which is the sum of subspaces \(\mathbf{X}_i (i \in I)\) , and let \(f\) be a continuous mapping of X into a topological space Y. For each \(i \in I\) , let \(f_i: \mathbf{X}_i \to \mathbf{Y}\) be the restriction of \(f\) to \(\mathbf{X}_i\) . Show that \(f\) is proper if and only if each of the \(f_i\) is proper and the family \((f(\mathbf{X}_i))_{i \in I}\) is locally finite {[}use Theorem 1 d){]}.

\begin{enumerate}
\def\labelenumi{\arabic{enumi})}
\setcounter{enumi}{6}
\item
  Let \(f: \mathbf{X} \to \mathbf{X}'\) be a proper mapping and let \(\mathbf{R}\) (resp. \(\mathbf{R}'\) ) be an equivalence relation on \(\mathbf{X}\) (resp. \(\mathbf{X}'\) ) with which \(f\) is compatible. Let \(\overline{f}: \mathbf{X} / \mathbf{R} \to \mathbf{X}' / \mathbf{R}'\) be the mapping of the quotient spaces induced by \(f\) . Suppose that (i) the image under \(f\) of every saturated set with respect to \(\mathbf{R}\) is saturated with respect to \(\mathbf{R}'\) , and (ii) \(\mathbf{R}'\) is open. Show that in these conditions \(\overline{f}\) is proper {[}use Theorem 1 c){]}.
\item
  Show that under the hypotheses of Example 3 of § 5, no. 2, the canonical mapping \(\mathbf{X} \to \mathbf{X} / \mathbf{R}\) is proper.
\end{enumerate}

¶ 9) a) Let X, Y be two topological spaces and \(f: \mathbf{X} \to \mathbf{Y}\) a surjective closed (but not necessarily continuous) mapping. Show that if \(\overline{f}^{1}(y)\) is quasi-compact for each \(y \in \mathbf{Y}\) , then \(\overline{f}^{1}(\mathbf{K})\) is quasi-compact for each quasi-compact subset K of Y (either prove the result directly or else use Exercise 7 of § 5 and Exercise 12 of § 9).

\begin{enumerate}
\def\labelenumi{\alph{enumi})}
\setcounter{enumi}{1}
\tightlist
\item
  Suppose in addition that X is regular. Show that if \(\overline{f}^{1}(y)\) is closed in X for each \(y \in Y\) , then \(\overline{f}^{1}(K)\) is closed in X for each quasi-compact subset K of Y (same methods). If also Y is locally compact then f is continuous.
\end{enumerate}

\begin{enumerate}
\def\labelenumi{\arabic{enumi})}
\setcounter{enumi}{9}
\tightlist
\item
  Let X, Y be two topological spaces. A correspondence (G, X, Y) between X and Y (Set Theory, Chapter II, § 3, no. 1) is said to be proper if the restriction of the projection \(\mathbf{pr}_2\) to the graph G is a proper mapping \(\mathbf{G} \to \mathbf{Y}\) .
\end{enumerate}

\begin{enumerate}
\def\labelenumi{\alph{enumi})}
\item
  Let \(f: \mathbf{X} \to \mathbf{Y}\) be a mapping. Show that \(f\) is continuous if and only if the correspondence \(\overline{f}^1\) between \(\mathbf{Y}\) and \(\mathbf{X}\) is proper, and that \(f\) is a proper mapping if and only if the correspondences \(\overline{f}^1\) (between \(\mathbf{Y}\) and \(\mathbf{X}\) ) and \(f\) (between \(\mathbf{X}\) and \(\mathbf{Y}\) ) are proper. Hence construct an example of a non-continuous mapping \(f\) such that the correspondence \(f\) is proper.
\item
  If \((\mathbf{G}, \mathbf{X}, \mathbf{Y})\) is a proper correspondence, show that \(\mathbf{G}(\mathbf{A})\) is a closed subset of Y whenever A is a closed subset of X.
\end{enumerate}

¶ 11) Let (G, X, Y) be a proper correspondence (Exercise 10).

\begin{enumerate}
\def\labelenumi{\alph{enumi})}
\tightlist
\item
  Show that if \(\mathfrak{B}\) is a filter base on \(X\) which consists of closed sets, then
\end{enumerate}

\[
\bigcap_ {\mathbf {M} \in \mathfrak {B}} \mathrm{G} (\mathbf {M}) = \mathrm{G} \left(\bigcap_ {\mathbf {M} \in \mathfrak {B}} \mathbf {M}\right).
\]

\begin{enumerate}
\def\labelenumi{\alph{enumi})}
\setcounter{enumi}{1}
\tightlist
\item
  For each \(y \in \mathbf{Y}\) and each neighbourhood \(\mathbf{V}\) of \(\overline{\mathbf{G}}(y)\) in \(\mathbf{X}\) , show that there is a neighbourhood \(\mathbf{W}\) of \(y\) in \(\mathbf{Y}\) such that \(\overline{\mathbf{G}}^{1}(\mathbf{W}) \subset \mathbf{V}\) (``continuity of the roots as functions of the parameters'') (use Exercise 6 of § 9).
\end{enumerate}

¶ 12) a) A correspondence (G, X, Y) is proper if and only if, for every topological space Z and every correspondence (E, Z, X) whose graph E is closed in Z × X, the graph G \(\circ\) E is closed in Z × Y. First to show that the condition is necessary, remark that G \(\circ\) E is the image of (E × Y) ∩ (Z × G) under \(\iota_{\mathbf{Z}} \times \mathrm{pr}_{2} : \mathbf{Z} \times \mathbf{G} \to \mathbf{Z} \times \mathbf{Y}\) . To show sufficiency, let \(\delta_{\mathbf{Y}}\) be the diagonal mapping \(\mathbf{Y} \to \mathbf{Y} \times \mathbf{Y}\) ; show that, for each subset \(\mathbf{F}\) of \((\mathbf{Z} \times \mathbf{Y}) \times \mathbf{X}\) , the inverse image of \(\mathbf{G} \circ \mathbf{F}\) under the mapping

\[
\iota_ {\mathbf {Z}} \times \delta_ {\mathbf {Y}}: \mathbf {Z} \times \mathbf {Y} \rightarrow \mathbf {Z} \times \mathbf {Y} \times \mathbf {Y}
\]

is the image, under the mapping \((y, z) \to (z, y)\) , of the set

\[
\operatorname{pr} _ {2 3} (\overline {{{\mathrm{F}}}} \cap (\mathrm{G} \times \mathrm{Z})),
\]

where \(\mathbf{F}\) is the image of \(\mathbf{F}\) under the mapping \((z, y, x) \to (x, y, z)\) . b) Deduce from a) that the composition of two proper correspondences is a proper correspondence.

\begin{enumerate}
\def\labelenumi{\alph{enumi})}
\setcounter{enumi}{2}
\tightlist
\item
  Let \((\mathbf{G},\mathbf{X},\mathbf{Y})\) be a proper correspondence. Show that if \(\mathbf{X}\) is Hausdorff then \(\mathbf{G}\) is closed in \(\mathbf{X} \times \mathbf{Y}\) .
\end{enumerate}

¶ 13) Let X be a locally compact space and Y any topological space. a) Show that the following four statements are equivalent:

\(\alpha)\) (G, X, Y) is a proper correspondence between X and Y.

\(\beta)\) For each Hausdorff space \(\mathbf{X}'\) which contains \(\mathbf{X}\) as a subspace, \(\mathbf{G}\) is closed in \(\mathbf{X}' \times \mathbf{Y}\) .

\(\gamma)\) G is closed in \(\mathbf{X} \times \mathbf{Y}\) , and for each \(y \in \mathbf{Y}\) there is a neighbourhood \(\mathbf{V}\) of \(y\) in \(\mathbf{Y}\) and a compact subset \(\mathbf{K}\) of \(\mathbf{X}\) such that \((\mathbf{X} - \mathbf{K}) \times \mathbf{V}\) does not meet \(\mathbf{G}\) .

δ) There is a compact space \(X'\) , containing X as a subspace, such that G is closed in \(X' \times Y\) .

{[}To prove that \(\alpha) \Longrightarrow \beta)\) use Exercise 12 \(a)\) and \(c)\) . To show that \(\delta) \Longrightarrow \alpha)\) , use Proposition 2 of no. 1 and Corollary 5 of Theorem 1 of no. 2.{]}

\begin{enumerate}
\def\labelenumi{\alph{enumi})}
\setcounter{enumi}{1}
\tightlist
\item
  Show that if \((\mathbf{G},\mathbf{X},\mathbf{Y})\) is proper, then:
\end{enumerate}

\(\zeta)\) For every compact subset \(\mathbf{L}\) of \(\mathbf{Y}\) , \(\overline{\mathbf{G}}(\mathbf{L})\) is a compact set.

Show also that if Y is locally compact and if G is closed in \(X \times Y\) and satisfies \(\zeta\) , then (G, X, Y) is proper. {[}Use the equivalence of \(\alpha\) and \(\gamma\) .{]} This last assertion is no longer valid when Y is Hausdorff but not locally compact (Exercise 4).

\begin{enumerate}
\def\labelenumi{\arabic{enumi})}
\setcounter{enumi}{13}
\item
  Let Y be a locally compact space, X a topological space. Show that a mapping \(f: X \rightarrow Y\) is continuous if and only if, for each compact space \(Y'\) which contains Y as a subspace, the graph of f is closed in \(X \times Y'\) . {[}Use Exercises 13 and 10 a){]}.
\item
  Let X be a Hausdorff space, A a closed subset of X, f a proper mapping of X --- A into a locally compact space Y, Y' the one-point compactification of Y, and \(\omega\) the point at infinity in Y'. Let g be the mapping of X into Y' which coincides with f on X --- A and maps each point of A to \(\omega\) . Show that g is continuous.
\end{enumerate}

¶ 16) Let X be a locally compact space, R an equivalence relation on X, C the graph of R in X × X. Show that C is closed in X × X if and only if, for each compact subset K of X, the saturation of K with respect to R is closed in X. (To prove necessity, use Corollary 5 of Theorem 1 of no. 2).

*17) On the locally compact space R, let S be the equivalence relation obtained by identifying all the points of Z. Show that S is closed and R/S Hausdorff but not locally compact.*

*18) Let X be the locally compact subspace {[}o, +∞{[} of R. Let S be the equivalence relation on X whose equivalence classes are the sets \{x, 1/x\} for o \textless{} x ≤ 1 and the set \{o\}. Show that S is open, that the graph of S is closed in X × X and that X/S is compact, but that S is not closed.*

¶ 19) Let X be a locally compact σ-compact space, R an equivalence relation on X whose graph is closed in X × X. Show that X/R is Hausdorff. {[}Let (Uₙ) be a covering of X by relatively compact open sets such that Uₙ ⊂ Uₙ₊₁. Let A, B be two disjoint closed sets which are saturated with respect to R. Define recursively two increasing sequences (Vₙ), (Wₙ) of closed sets saturated with respect to R, such that Vₙ ∩ Wₙ = ∅ and such that Vₙ ∩ Uₙ (resp. Wₙ ∩ Uₙ) is a neighbourhood of A ∩ Uₙ (resp. B ∩ Uₙ) in the compact subspace Uₙ. Use Proposition 8 of no. 3 and Exercise 16, as well as Proposition 2 of § 9, no. 2{]}. (*)

¶ 20) a) Let X be a non-empty compact space and R a Hausdorff equivalence relation on X. Show that X/R is closed in the space \(\mathfrak{P}_{0}(X)\) with the topology \(\mathcal{G}_{\Theta}\) (§ 8, Exercise 12) if and only if the relation R is open. {[}Use Exercise 13 of § 9 and Exercise 16 of § 3 to show that if X/R is closed in \(\mathfrak{P}_{0}(X)\) in the topology \(\mathcal{G}_{\Theta}\) , then the quotient topology on X/R coincides with that induced by \(\mathcal{G}_{\Theta}\) , and then apply Exercise 7 of § 5. Conversely, if R is open, use Exercise 29 of § 8 to show that X/R is closed with respect to \(\mathcal{G}_{\Theta}\) {]}.

*b) Let X be the locally compact subspace of the real plane \(\mathbb{R}^2\) consisting of the points \((\mathtt{I}/n, y)\) , where \(n\) is an integer \(>0\) and \(-\mathtt{I} \leqslant y \leqslant \mathtt{I}\) , and the points \((\mathtt{o}, y)\) where either \(-\mathtt{I} \leqslant y < 0\) or \(0 < y \leqslant \mathtt{I}\) or \(y = 2\) . Let \(p\) be the restriction of \(\mathsf{pr}_1\) to X, and R the equivalence relation associated with \(p\) . Show that R is neither open nor closed, but that X/R is compact and closed in \(\mathfrak{P}_0(X)\) in the topology \(\mathcal{C}_{\Theta}\) . * ¶ 21) (*) a) Let X be an accessible space which is locally quasi-compact (§ 9, Exercise 29). Let \(\mathbf{X}_0\) be the discrete space which has the same underlying set as X, and let \(\tilde{\mathbf{X}}_0\) be the (compact) ultrafilter space of \(\mathbf{X}_0\) (the discrete space \(\mathbf{X}_0\) being identified with an open subspace of \(\tilde{\mathbf{X}}_0\) ) (§ 9, Exercise 26). On \(\tilde{\mathbf{X}}_0\) , let R {[}or R(X){]} be the following equivalence relation between ultrafilters U, B: ``the sets of limit points of U and B in X are the same''. Show that R is Hausdorff {[}use the description of the topology of \(\tilde{\mathbf{X}}_0\) given in § 9, Exercise 25 b), Exercise 7 b) of § 7 and axiom (C'){]}. The compact space \(\mathbf{X}^c = \tilde{\mathbf{X}}_0 / \mathbf{R}\) is in one-to-one correspondence with the set of primitive subsets of X (§ 7, Exercise 7). If X is Hausdorff (and therefore locally compact) then \(\mathbf{X}^c\) can be identified with X if X is compact, and with the one-point compactification of X if X is not compact.

\begin{enumerate}
\def\labelenumi{\alph{enumi})}
\setcounter{enumi}{1}
\item
  Let Y be a closed subspace of X. The ultrafilter space \(\tilde{Y}_{0}\) can be identified with the closure of \(Y_{0}\) in \(\tilde{X}_{0}\) , and \(Y^{c}\) with the canonical image of \(\tilde{Y}_{0}\) in \(X^{c}\) .
\item
  Show that the restriction to \(X_{0}\) of the canonical mapping \(\tilde{X}_{0} \rightarrow X^{c}\) is a bijection of \(X_{0}\) onto a subspace \(X_{*}^{c}\) of \(X^{c}\) , that the canonical mapping \(f: X_{*}^{c} \rightarrow X\) is a continuous bijection, and that the image under this mapping of any compact subset of \(X_{*}^{c}\) is a closed compact subset of X. {[}To show that f is continuous, use b); also remark that the inverse image in \(\tilde{X}_{0}\) of a point of \(X_{*}^{c}\) consists of all the ultrafilters on X which have a single given limit point, and that if U is an ultrafilter on X, then its limit point in \(\tilde{X}_{0}\) , when U is regarded as an ultrafilter base on \(\tilde{X}_{0}\) , is U regarded as a point of \(\tilde{X}_{0}\) {]}.
\end{enumerate}

*d) Let \(\mathbf{Z}\) be the set which is the sum of a countable number of spaces \(\mathbf{Y}_n\) each identical with the space \(\mathbf{Y}\) defined in Exercise 7 of §8 and a set \(\{\omega\}\) consisting of a single point. Define a topology on \(\mathbf{Z}\) by taking as a fundamental system of neighbourhoods of a point \(z \in \mathbf{Y}_n\) its neighbourhood filter in the space \(\mathbf{Y}_n\) , and as fundamental system of neighbourhoods of \(\omega\) the set of subsets \(\mathbf{V}_n\) of \(\mathbf{Y}\) , where \(\mathbf{V}_n\) is the union of \(\{\omega\}\) and the \(\mathbf{Y}_m\) where \(m \geqslant n\) . The space \(\mathbf{Z}\) so defined is accessible, quasi-compact and locally quasi-compact; but the subspace \(Z_*^c\) of \(Z^c\) is not locally compact.*

\begin{enumerate}
\def\labelenumi{\alph{enumi})}
\setcounter{enumi}{4}
\tightlist
\item
  Let Y be another locally quasi-compact accessible space, and let u be a continuous mapping of X into Y, such that the inverse image under u of each quasi-compact subset of Y is quasi-compact. Regarding u as a mapping of \(X_{*}^{c}\) into \(Y_{*}^{c}\) , show that u is continuous and can be (uniquely) extended to a continuous mapping \(\mathbf{X}^c\to \mathbf{Y}^c\) . {[}First extend \(u\) to a continuous mapping of \(\tilde{\mathbf{X}}_0\) into \(\tilde{\mathbf{Y}}_0\) and show that the extension is compatible with the equivalence relations \(\mathbf{R}(\mathbf{X}),\mathbf{R}(\mathbf{Y})\) {]}.
\end{enumerate}

\(^{*}f)\) Let X be the quasi-compact, not locally quasi-compact, accessible space defined in Exercise 29 c) of § 9. Show that in the corresponding space \(\tilde{X}_{0}\) the equivalence relation R(X) is not Hausdorff.*

